\documentclass[11pt]{article}
\usepackage[margin=1in]{geometry}
\usepackage{amsmath,amssymb,amsthm,mathtools}
\usepackage{array,booktabs,longtable}
\usepackage[hidelinks]{hyperref}

\newtheorem{theorem}{Theorem}[section]
\newtheorem{lemma}[theorem]{Lemma}
\newtheorem{proposition}[theorem]{Proposition}
\newcommand{\spE}{s^+}
\newcommand{\smE}{s^-}
\newcommand{\sig}{\sigma}
\newcommand{\one}{\mathbf 1}

\title{Positive Square Energy of Every Octacyclic Cactus}
\author{}
\date{26 July 2026}

\begin{document}
\maketitle

\begin{abstract}
Let $G$ be a connected octacyclic cactus on $n$ vertices.  We prove
$s^+(G)>n$.  Sharp block-additive doubly nonnegative optimization reduces the
problem to $T^7Q$ and $T^6PP$, where $T=C_3$, $P=C_5$, and $Q$ is an arbitrary
cycle.  For disconnected shared-cut graphs, exact colored-partition censuses
give $44=42+2$ and $76=70+6$ proper rows.  The structural $T^7Q$ endpoint is
closed by a standalone packing-one Sachs lemma whose hypothesis is checked on
a common-cut triangle bouquet.  The entry-locked $T^6P\mid P$ endpoint is
closed by the corrected strict-last-bridge certificate $877=861+16$, including
all sixteen marked profiles and the distinction between a shared-cut $TTP$
packet and a common-cut bouquet.  For one shared-cut cluster, exact incidence
censuses leave only the common-cut $T^7Q$ bouquet and six of the $2116$ types
in $T^6PP$; the common-cut Schur--Sachs theorem and six explicit router
resolutions close them.  Every partition is induced and assigns every
connector remnant, router interval, shared cut, and attached tree exactly once.
The proof expressly does not use the retracted all-rank rooted hostile-cycle
guard or the superseded uncut $877=868+9$ certificate.
\end{abstract}

\section{Definitions and validity boundary}

All graphs are finite, simple, and undirected.  A cactus is a connected graph
whose blocks are bridges or cycles.  Its cyclomatic number equals its number of
cyclic blocks.  If the adjacency eigenvalues of $G$ are
$\lambda_1,\ldots,\lambda_n$, put
\[
 \spE(G)=\sum_{\lambda_i>0}\lambda_i^2,
 \qquad \smE(G)=\sum_{\lambda_i<0}\lambda_i^2,
 \qquad \sig(G)=\spE(G)-|V(G)|.
\]
For every vertex partition $V(G)=V_1\mathbin{\dot\cup}\cdots
\mathbin{\dot\cup}V_r$, induced-subgraph superadditivity gives
\begin{equation}\label{eq:super}
 \spE(G)\geq\sum_i\spE(G[V_i]),
 \qquad \sig(G)\geq\sum_i\sig(G[V_i]).
\end{equation}
This is the theorem of Akbari, Kumar, Mohar, and Pragada~\cite{AKMP}.
Akbari, Kumar, Mohar, Pragada, and Zhang conjectured that
$s^+(G)\geq |V(G)|$ for every connected graph with at least
$|V(G)|+1$ edges~\cite[Conjecture~1.2]{AKMPZ}.  Since an octacyclic graph has
$|E(G)|=|V(G)|+7$, Theorem~\ref{thm:main} proves their conjecture, strictly,
for every connected octacyclic cactus.

Write $T=C_3$, $P=C_5$, and
\[
 \delta=\sec(\pi/5)-1=\sqrt5-2<\frac14,
 \qquad \delta_q=\sec(\pi/q)-1<1.
\]
The symbol $Q$ denotes $C_q$.  It is called hostile only when
$q\equiv1\pmod4$; otherwise its unicyclic packet is nonnegative.

This manuscript uses only the corrected provenance summarized in
\cite{Synthesis}: the strict-last-bridge $877=861+16$ certificate
\cite{LastBridge}, the standalone packing-one hostile-cycle lemma
\cite{PackingOne}, the exact fully shared census and six replacements
\cite{SharedCensus,Six}, and the scalar common-cut theorem
\cite{CommonCut}.  We explicitly reject and do not use the retracted all-rank
rooted hostile-cycle guard.  In particular, we never assert that a maximum
triangle packing makes every nearest-cycle Voronoi territory packing one.
We also do not use the superseded uncut $877=868+9$ (E1--E9) computation; its
connector ownership was defective.  Every occurrence of ``packing one'' below
is the standalone lemma with its hypothesis verified directly.

There are two intentionally different terms.  A \emph{shared-cut $TTP$}
packet is a connected packet in which its two triangles share a cut; its
pentagon belongs to the packet but need not contain that cut.  A
\emph{common-cut $T^kP$} or \emph{common-cut $T^kPP$} packet has one pivot
contained in every named cycle.  The former is never silently treated as the
latter.

\section{Packet bounds and territory operations}

We use the established lower-rank and packet estimates
\begin{align}
 \sig(H)&\geq0 &&\text{if $H$ is bicyclic or tricyclic},\label{eq:low0}\\
 \sig(H)&>0 &&\text{if $H$ has rank $4,5,6$, or $7$},\label{eq:lowpos}\\
 \sig(P)&\geq-\delta,& \sig(Q)&\geq-\delta_q
   &&\text{for hostile $Q$},\label{eq:hostile}\\
 \sig(TP)&>1-\delta,\label{eq:tp}\\
 \sig(PP)&>0 &&\text{for one shared-cut cluster},\label{eq:pp}\\
 \sig(TPP)&>\frac32 &&\text{for one shared-cut cluster},\label{eq:tpp}\\
 \sig(TTP)&>2-\delta
   &&\text{when the two triangles share a cut}.\label{eq:ttp}
\end{align}
All permit arbitrary finite trees at arbitrary core vertices; $TP$ permits an
arbitrary bridge connector.  The hostile unicyclic bound is
\cite[Lemma~1.1(2)]{Tricyclic}, the mixed $TP$ bound is
\cite[Theorem~3.1]{Bicyclic}, and the complete bicyclic and tricyclic bounds are
\cite[Theorem~5.1]{Bicyclic} and \cite[Theorem~5.1]{Tricyclic}.  The shared-cut
$PP$ bound is \cite[Theorem~4.1]{Bicyclic}; the shared-cluster $TPP$ and
shared-cut $TTP$ bounds are \cite[Proposition~4.2]{Tricyclic} and
\cite[Proposition~3.1]{Tricyclic}.  The strict rank-four input is the repaired
tetracyclic theorem \cite[Theorem~5.1]{Tetracyclic}: its proof
uses the exact twenty-core $C_{3355}$ audit and common-root phase repair, not
the false raw $C_{3335}$ coefficient certificate.  The strict rank-five,
rank-six, and rank-seven inputs are \cite[Theorem~8.1]{Pentacyclic},
\cite[Theorem~9.1]{Hexacyclic}, and \cite[Theorem~9.1]{Heptacyclic}.
If $A_r$ is one connected shared-cut cluster of $r$ triangles, successive
incidence-leaf openings from the corrected rank-four matching injection give
\begin{equation}\label{eq:triangle}
 \sig(A_r)>b_r,
 \qquad (b_1,\ldots,b_7)=(0,1,2,3,2,1,0).
\end{equation}
In particular, $\sig(A_4)>3$ is not obtained from a multiplicity-blind Sturm
count.  The complete recurrence, including arbitrary attached trees, is
\cite[Lemma~3.1]{Heptacyclic}; its rank-four exceptional base is the exact
matching-injection result \cite[Proposition~4.11]{PackingTwo}.

Two analytic results are essential.

\begin{lemma}[Standalone packing-one hostile packet]\label{lem:packing-one}
Let a connected cactus have one hostile $Q=C_q$ and $a\geq1$ triangles, no
two of which are vertex-disjoint.  The $Q$ may meet the triangular part at a
root or be joined to its root by an arbitrary positive path.  Arbitrary trees
may occur everywhere, including on that path.  Then
\[
 \spE(H)-\smE(H)>-2\delta_q,
 \qquad \sig(H)>a-\delta_q.
\]
\end{lemma}

\begin{proof}
Let $S=V(Q)$ and let the spine be the cyclic blocks and their minimal joining
paths.  Eliminate every component off the spine by the exact signless matching
message
\[
 M_{u\to v}(t)=\frac{Z_{T_{u\to v}}(t)}{Z_{T_{u\to v}-u}(t)}
 =t+\sum_wM_{w\to u}(t)^{-1}\geq t.
\]
This replaces each spine activity by $\alpha_v=t+y_v\geq t$ and removes a
common positive factor.  Write $Z_F(\alpha)$ for the resulting signless
matching partition.  The grouped Sachs expansion on the imaginary axis is
\[
 \Psi_H(t)=R+2i(B-A),
\]
where, after normalization,
\[
 B=Z_{H-S}(\alpha),\quad A=\sum_T Z_{H-V(T)}(\alpha),\quad
 R=Z_H(\alpha)+4\!\sum_{T\cap Q=\varnothing}
 Z_{H-(S\cup V(T))}(\alpha)>0.
\]
Packing one is used exactly here: no Sachs collection contains two triangles;
the last sum is precisely the admissible $Q$--triangle terms.  If $Z_q(t)$ is
the bare matching partition of $Q$, split the matchings in $Z_H(\alpha)$ by
whether they use a boundary edge of $S$.  Then
\[
 Z_H(\alpha)=Z_Q(\alpha|_S)B+E,\quad E\geq0,
 \qquad Z_Q(\alpha|_S)=Z_q(t)+L,\quad L\geq0,
\]
and hence
\[
 R-Z_q(t)(B-A)=E+LB+Z_q(t)A
 +4\!\sum_{T\cap Q=\varnothing}Z_{H-(S\cup V(T))}(\alpha)>0.
\]
If $B-A\leq0$ the argument is nonpositive; otherwise division by
$RZ_q(t)>0$ gives the same conclusion.  Thus the continuous argument of
$\Psi_H(t)$ is strictly smaller than $\arctan(2/Z_q(t))$, the argument for
isolated $Q$.  The signed Coulson identity
\[
 \spE(F)-\smE(F)=-\frac4\pi\int_0^\infty t\,\operatorname{Arg}\Psi_F(t)\,dt
\]
and $\spE(C_q)-\smE(C_q)=-2\delta_q$ give the first inequality.
Since $|E(H)|=|V(H)|+a$ and $\spE+\smE=2|E|$, the second follows.  This is the
standalone grouped-Sachs proof of~\cite{PackingOne}; it has no all-rank or
Voronoi premise.
\end{proof}

\begin{lemma}[Common-cut Schur--Sachs theorem]\label{lem:common-cut}
For a one-vertex bouquet, uniformly over arbitrary attached trees,
\begin{align}
 \sig(T^kQ)&>k-\delta_q &&(q\equiv1\pmod4),\label{eq:ccq}\\
 \sig(T^kP)&>k-\delta &&(k\geq1),\label{eq:ccp}\\
 \sig(T^kQ)&>k &&(q\text{ even or }q\equiv3\pmod4),\label{eq:ccqgood}\\
 \sig(T^kPP)&>k+1-\frac4{3\sqrt{13}}.\label{eq:ccpp}
\end{align}
\end{lemma}

\begin{proof}
After exact Schur elimination of attached trees, each core activity is
$a_v(t)=t+y_v(t)$ with $y_v\geq0$.  If $A_j$ is the matching partition of
the private path in lobe $j$ and $B_j$ is the contribution obtained by
matching the pivot into that lobe, the normalized determinant is the scalar
expression
\[
 W(t)=a_x+\sum_j\frac{B_j+\omega_{\ell_j}}{A_j},
 \qquad \omega_\ell=-2i^{-\ell}.
\]
For $T^kQ$, put
\[
 X=a_x+B_Q/A_Q+\sum_{j=1}^kB_j/A_j>0.
\]
For hostile $Q$ the imaginary part is
$2(1/A_Q-\sum_j1/A_j)$, while $A_QX>Z_q(t)$ because $k\geq1$, $B_j>0$, and
all activities are at least $t$.  If that imaginary part is nonpositive the
phase is nonpositive; otherwise
\[
 \frac{2(1/A_Q-\sum_j1/A_j)}X<\frac2{A_QX}<\frac2{Z_q(t)}.
\]
Thus its argument is strictly below the isolated hostile-$Q$ argument, and is
negative for the other congruence classes.  Coulson integration gives
\eqref{eq:ccq} and \eqref{eq:ccqgood}.  Taking $Q=P=C_5$ gives the explicit
specialization \eqref{eq:ccp}; this is the common-cut $T^kP$ estimate used in
L1--L12 and U2--U3.  For $T^kPP$, write
$X_{PP}=a_x+B_1/A_1+B_2/A_2$ and $Y_{PP}=2/A_1+2/A_2$.
Adding triangles strictly increases the positive real part and decreases the
imaginary part, so its phase is strictly below that of the two-pentagon core.
The exact positive-coefficient inequality
\[
 2R\geq t(t^4+7t^2+9)(A_1+A_2)
\]
implies
$\Theta_{PP}(t)\leq4/[t(t^4+7t^2+9)]$.  Since
\[
 \int_0^\infty\frac4{t^4+7t^2+9}\,dt=\frac{2\pi}{3\sqrt{13}},
\]
the signed energy difference exceeds $-8/(3\sqrt{13})$, proving
\eqref{eq:ccpp}.  This scalar theorem requires the actual common pivot and is
not applied to multi-pivot shared-cut packets.
\end{proof}

Join cyclic blocks when they share a cut and call the components
\emph{shared-cut clusters}.  Contract each cluster in the block-cut tree and
retain the actual bridge paths between clusters; the resulting reduced cluster
graph is a tree.

\begin{lemma}[Whole-cluster realization]\label{lem:whole-cluster}
Let $S_1,\ldots,S_r$ be pairwise disjoint connected subtrees of the reduced
cluster tree and suppose every marked shared-cut cluster belongs to exactly one
$S_i$.  Then $V(G)$ has a partition into connected induced territories
$G_1,\ldots,G_r$ such that $G_i$ retains, whole, exactly the cyclic blocks in
clusters assigned to $S_i$.  Every separation is at an actual bridge, and
every connector remnant, unmarked Steiner branch, connector entry, and
off-hull tree has exactly one owner.
\end{lemma}

\begin{proof}
Contract each $S_i$ to a labelled root and assign every remaining reduced-tree
vertex to a nearest root, breaking ties by a fixed label order.  Each label
class in a tree is connected.  An edge between classes expands in the
block-cut tree to a path of bridge blocks; cut one actual bridge on that path
and give the two remnants to the adjacent owners.  At a Steiner branch, give
the branch vertex and one initial segment to one owner and cut an actual bridge
on every other incident segment.  A component off the cyclic hull has a unique
hull attachment, or the block-cut graph would contain a cycle, and is assigned
whole to that attachment's owner.  The resulting territories are connected,
induced, disjoint, exhaustive, and retain the claimed whole clusters.
\end{proof}

\begin{lemma}[Accepted one-cycle interval realization]\label{lem:ownership}
Let one cycle $C$ in a shared-cut cluster be selected by a finite certificate.
Group its incidence branches and external connector entries by their vertex on
$C$; coincident marks form one group.  If the distinct groups can be assigned
to nonempty proper consecutive intervals of $C$, then each interval together
with all branches and entries assigned to it is a connected induced territory,
the territories partition everything assigned at $C$, and $C$ is retained by
none.  A later accepted split may refine one territory.  Every coincident entry
at a shared cut follows the unique interval owning that cut.
\end{lemma}

\begin{proof}
List the distinct marked vertices cyclically.  Cut one cycle edge in each
chosen gap and assign every unmarked gap vertex to one adjacent interval.  Each
interval is a proper path.  The cycle--cut incidence graph is a tree, so every
branch of $I-C$ meets $C$ at exactly its marked vertex and follows exactly one
interval.  Entries coincident with that vertex follow the same owner and do
not demand an empty extra interval.  A refinement partitions only one already
induced vertex set; hence inducedness, disjointness, exhaustiveness, and unique
ownership persist sequentially.
\end{proof}

\section{Exact DNN reduction}

For a connected edge-containing graph set
\[
 \kappa(G)=\sup_{\substack{M\succeq0,\ M\geq0\\\one^TM\one>0}}
 \frac{4\bigl(\sum_{uv\in E(G)}\sqrt{M_{uv}}\bigr)^2}{\one^TM\one}.
\]
The sharp cactus DNN theorem~\cite{DNN,LTZ} says that for bridge-block count
$b$ and cycle lengths $\ell_1,\ldots,\ell_8$,
\[
 \kappa(G)=b+\sum_i\kappa(C_{\ell_i}),\qquad
 \kappa(C_\ell)=\begin{cases}
  \ell,&\ell\text{ even},\\
  \ell\sec^2(\pi/(2\ell)),&\ell\text{ odd},
 \end{cases}
 \qquad \smE(G)\leq\kappa(G).
\]
Put $\epsilon_\ell=0$ for even $\ell$ and
$\epsilon_\ell=\ell\tan^2(\pi/(2\ell))$ for odd $\ell$.  Since an
octacyclic connected cactus has $|E|=n+7$, block counting gives
\begin{equation}\label{eq:dnn}
 \sig(G)\geq7-\sum_{i=1}^8\epsilon_{\ell_i}.
\end{equation}
The odd sequence decreases, $\epsilon_3=1$, and with
$a=\epsilon_5=5-2\sqrt5$ the exact comparisons are
\[
 3a<2,\qquad2a>1,\qquad\epsilon_5+\epsilon_7<1.
\]
Hence at most five triangles gives
$\sum\epsilon_i\leq5+3a<7$.  With exactly six triangles the remaining pair
reaches one only for $P,P$.  With at least seven triangles the multiset is
$T^7Q$.  Therefore~\eqref{eq:dnn} is strict outside exactly
\begin{equation}\label{eq:residuals}
 T^7Q=\{3,3,3,3,3,3,3,q\},\qquad
 T^6PP=\{3,3,3,3,3,3,5,5\}.
\end{equation}

\section{Disconnected shared-cut graph}

The census uses exact rational arithmetic from Python's \texttt{fractions}
module and gives
\begin{center}
\begin{tabular}{lrrr}
\toprule
residual & all partitions & proper partitions & direct packet rows\\
\midrule
$T^7Q$  &45&44&42\\
$T^6PP$ &77&76&70\\
\bottomrule
\end{tabular}
\end{center}
A direct row has a positive packet sum, with strictness tracked symbolically.
The only structural rows are
\begin{align}
 &Q|T|T|T|T|T|T|T,\qquad Q|T^7,\label{eq:t7rows}\\
 &P|P|T|T|T|T|T|T,\quad P|T|T|T|T|T^2P,\quad
 P|T|T|T|T^3P,\nonumber\\
 &P|T|T|T^4P,\quad P|T|T^5P,\quad P|T^6P.\label{eq:t6rows}
\end{align}

\subsection{The two $T^7Q$ rows}

In the all-singleton row, a singleton triangle is a leaf of the finite reduced
tree; cutting its first bridge gives a strict triangle and a connected
heptacyclic complement.  Consider $T^7|Q$.  If the triangle incidence tree is
not a bouquet, split an internal triangle at its cyclic marks and, when needed,
at the private connector-entry mark.  If the $Q$ branch retains $r$ triangles,
$r\geq3$ is positive by lower rank, $r=2$ is nonnegative with another strict
triangle branch, and $r=1$ is a positive $TQ$ packet.  If $r=0$, the entry uses
one of at most three marks, so six triangles occupy at most two other branches;
one contains at least three triangles and has surplus $>2$, absorbing the
isolated hostile loss $>-1$.

If the seven triangles form a bouquet but the connector enters a private
triangle vertex, splitting there gives $A_6+Q$ and surplus
$>1-\delta_q>0$.  The only locked structure is
\begin{equation}\label{eq:g7q}
 \text{seven triangles sharing $x$, joined from $x$ to $Q$ by an arbitrary
 positive connector.}
\end{equation}
Its triangles are pairwise intersecting, so their vertex-packing number is
exactly one.  Lemma~\ref{lem:packing-one}, with $a=7$, applies directly and
gives
\[
 \sig(G)>7-\delta_q>0.
\]
This verifies the standalone lemma's hypothesis on the packet itself; it is
not the retracted all-rank guard.  Nonhostile $Q$ also follows from
$A_7$ strictness and a bridge partition.  Thus all $44=42+2$ proper $T^7Q$
rows are proved.

\subsection{The $T^6PP$ rows before the locked endpoint}

For the first five rows of~\eqref{eq:t6rows}, a singleton-triangle leaf gives
$T$ and a heptacyclic complement.  If none is a leaf, the two nontriangle
endpoint clusters are the leaves of their path.  Pair the singleton triangle
nearest the singleton pentagon with that pentagon, obtaining a positive $TP$
terminal packet and a strict hexacyclic complement.  The separate row
$P_0|T^6|P_1$ has
\[
 \sig(G)>\sig(A_6)-2\delta>1-2\delta=5-2\sqrt5>0.
\]
Only $T^6P_0|P_1$ remains.  Splitting an internal $P_0$ into proper intervals
gives a mixed lower-rank entry branch and another strict triangular branch.  If
$P_0$ is an incidence leaf but the connector enters a private point of $P_0$,
split between entry and cyclic cut to obtain $A_6+P_1>1-\delta$.  The remaining
entry-locked family is denoted $G6PP$.

\subsection{\texorpdfstring{Strict-last-bridge $G6PP$: $877=861+16$}{Strict-last-bridge G6PP: 877=861+16}}

Cut the last actual bridge before remote $P_1$.  Its whole side is pentagonal
unicyclic and contributes at least $-\delta$.  The entire connector remnant on
the $T^6P_0$ side stays there as a tree rooted at its marked first cyclic-hull
entry.  Nothing is moved across the cut.  The color-preserving census has
\[
 226\text{ unrooted }T^6P_0\text{ incidences},\quad
 111\text{ with $P_0$ a leaf},\quad
 877\text{ marked root classes}.
\]
This is the exact root universe.  It starts with all $226$ color-preserving
bipartite cycle--cut trees on colors $T^6P_0$, restricts to the $111$ trees in
which $P_0$ has incidence degree one, and marks the first cyclic-hull entry
modulo color-preserving automorphism.  A mark is either any shared cut,
including the unique cut of $P_0$, or a private triangle vertex.  A triangle of
incidence degree $d$ contributes $3-d$ labelled private positions before orbit
quotienting.  Private vertices of $P_0$ are absent because the entry-locked
reduction has already disposed of that case.  Thus every possible locked
connector entry, and no unlocked entry, occurs among the $877$ classes.
The conservative one-router test proves $861$ and leaves exactly $16$, by cut
number $c=1,\ldots,6$:
\[
 (2,5,5,4,0,0).
\]
Here are the exact sixteen marked profiles.  Labels $0,\ldots,5$ are triangles,
$6=P_0$; $R$ marks a cut root and $TR$ a private triangle root.
\begin{center}
\small
\begin{longtable}{ccl}
\toprule
class&$c$&canonical rooted code\\
\midrule
\endfirsthead
\toprule
class&$c$&canonical rooted code (continued)\\
\midrule
\endhead
\midrule
\multicolumn{3}{r}{Continued on next page}\\
\endfoot
\bottomrule
\endlastfoot
L1&1&\texttt{R(P()T()T()T()T()T()T())}\\
L2&1&\texttt{X(P()T()T()T()T()T()TR())}\\
L3&2&\texttt{T(R(P()T()T()T()T())X(T()))}\\
L4&2&\texttt{T(R(T())X(P()T()T()T()T()))}\\
L5&2&\texttt{T(X(P()T()T()T()T())X(TR()))}\\
L6&2&\texttt{T(X(P()T()T()T()TR())X(T()))}\\
L7&2&\texttt{TR(X(P()T()T()T()T())X(T()))}\\
L8&3&\texttt{R(P()T()T()T(X(T()))T(X(T())))}\\
L9&3&\texttt{X(P()T()T()T(R(T()))T(X(T())))}\\
L10&3&\texttt{X(P()T()T()T(X(T()))T(X(TR())))}\\
L11&3&\texttt{X(P()T()T()T(X(T()))TR(X(T())))}\\
L12&3&\texttt{X(P()T()T(X(T()))T(X(T()))TR())}\\
L13&4&\texttt{R(P()T(X(T()))T(X(T()))T(X(T())))}\\
L14&4&\texttt{X(P()T(R(T()))T(X(T()))T(X(T())))}\\
L15&4&\texttt{X(P()T(X(T()))T(X(T()))T(X(TR())))}\\
L16&4&\texttt{X(P()T(X(T()))T(X(T()))TR(X(T())))}\\
\end{longtable}
\end{center}
The independent replacement verifier proves the following exhaustive ledgers.
$E$ is an acyclic private-entry interval, charged conservatively by
$\sig(E)\geq-1$.
\begin{center}
\begin{tabular}{lll}
\toprule
classes & clustered-side packets & total with $P_1$\\
\midrule
L1--L2 & common-cut $T^6P_0$ &$>6-2\delta$\\
L3--L6 & $T+$ common-cut $T^4P_0$ &$>4-2\delta$\\
L7 & $T+$ common-cut $T^4P_0+E$ &$>3-2\delta$\\
L8--L10,L12 & $T+T+$ common-cut $T^2P_0$ &$>2-2\delta$\\
L11 & $T+T+$ common-cut $T^2P_0+E$ &$>1-2\delta$\\
L13--L16 & $T+T+$ shared-cut $TTP_0$ &$>2-2\delta$\\
\bottomrule
\end{tabular}
\end{center}
For L13--L16 the two retained triangles share an actual cut, but $P_0$ need
not contain it; this is precisely the shared-cut $TTP$ distinction.  Its bound
is $>2-\delta$, not an invocation of the common-pivot theorem.  The weakest
case is L11:
\[
 0+0+(2-\delta)-\delta-1=1-2\delta=5-2\sqrt5>0.
\]
Binary routers are divided into a marked singleton and complementary edge;
saturated routers have three forced singleton intervals.  Every retained cut,
entry remnant, and attached tree has one owner.  More explicitly, L3--L7 split
$T_0$ and retain $T_2$ plus $T_1T_3T_4T_5P_0$; L8--L12 split $T_0$ first and
then $T_1$ inside the hub-owning territory, retaining $T_2,T_5,T_3T_4P_0$;
L13 splits the same routers and retains $T_2,T_4,T_3T_5P_0$; L14--L16 split
$T_1$ and then $T_3$, retaining $T_4,T_5,T_0T_2P_0$.  At each first split the
branches of the deleted router are disjoint components; the second split is
performed only inside the unique component owning the retained hub.  Every
retained cut is incident only with cycles in one final packet, while a private
entry on a deleted router becomes the separate acyclic $E$ owner exactly in
L7 and L11.  Lemma~\ref{lem:ownership} therefore proves sequential ownership,
not merely simultaneous ownership in the incidence code.  Thus all
$877=861+16$ classes
close, and all $76=70+6$ disconnected $T^6PP$ rows are proved.

\section{\texorpdfstring{One fully shared $T^7Q$ cluster}{One fully shared T7Q cluster}}

The exact color-preserving incidence census, by cut number $c=1,\ldots,7$, is
\begin{center}
\begin{tabular}{lrrrrrrr|r}
\toprule
$Q$ capacity&1&2&3&4&5&6&7&total\\
\midrule
$q=3$&1&9&49&142&236&191&60&688\\
$q=4$&1&9&49&145&243&202&66&715\\
$q=5$&1&9&49&145&245&205&69&723\\
$q=6$&1&9&49&145&245&206&70&725\\
$q\geq7$&1&9&49&145&245&206&71&726\\
\bottomrule
\end{tabular}
\end{center}
For every nonbouquet type, the exact ordinary-split ledger finds a legal
cycle split into proper consecutive intervals and a positive packet sum.  Its
branches use only $A_r$, isolated $Q$, $TQ$, $TTQ$, and rank four through seven
packets, with all retained shared-cut hypotheses checked.  The unique exception
in every capacity regime is the one-cut bouquet in which all seven triangles
and $Q$ share one vertex.  Lemma~\ref{lem:common-cut} gives there
\[
 \sig(T^7Q)>7-\delta_q>6\quad\text{for hostile $Q$},
 \qquad \sig(T^7Q)>7\quad\text{otherwise}.
\]
The theorem is used only in its actual common-pivot scope.  Hence every fully
shared $T^7Q$ cactus is proved.

\section{\texorpdfstring{One fully shared $T^6PP$ cluster: $2116=2110+6$}{One fully shared T6PP cluster: 2116=2110+6}}

The exact color-preserving census is
\begin{center}
\begin{tabular}{lrrrrrrr|r}
\toprule
&$c=1$&$c=2$&$c=3$&$c=4$&$c=5$&$c=6$&$c=7$&total\\
\midrule
all&1&14&106&377&728&657&233&2116\\
ordinary-split safe&0&13&104&376&727&657&233&2110\\
exceptions&1&1&2&1&1&0&0&6\\
\bottomrule
\end{tabular}
\end{center}
Each of the $2110$ accepted types has an exact positive packet ledger and a
legal proper-interval split.  The six remaining color-preserving types admit
the following complete replacements:
\begin{center}
\begin{tabular}{cll}
\toprule
code&replacement packetization&certified surplus\\
\midrule
U1&common-cut $T^6PP$&$>7-4/(3\sqrt{13})$\\
U2&$P+$ common-cut $T^5P$&$>5-2\delta$\\
U3&$P+T+$ common-cut $T^4P$&$>4-2\delta$\\
U4&$P+P+A_4$&$>3-2\delta$\\
U5&$P+P+T+A_3$&$>2-2\delta$\\
U6&$P+P+T+T+A_2$&$>1-2\delta$\\
\bottomrule
\end{tabular}
\end{center}
For completeness, the router realizations are as follows.  U1 is exactly the
common-cut $(6,2)$ bouquet and uses~\eqref{eq:ccpp} without splitting its
pivot.  In U2, a binary triangle router separates one pentagon and leaves an
actual common-cut $T^5P$.  In U3, one saturated router separates $P$, $T$, and
an actual common-cut $T^4P$.  In U4, two binary routers isolate the two
pentagons and leave $A_4$.  In U5, a saturated router and then a binary router
leave $P+P+T+A_3$.  In U6, two successive saturated-router splits leave
$P+P+T+T+A_2$.

At a binary router, the marked pentagon-side vertex is a singleton and the
other two router vertices form the complementary interval.  At a saturated
router all three marks occupy distinct triangle vertices, so the singleton
intervals are forced.  Each pentagon's unique cut remains with it; the common
hub has one owner; a second split only refines that owner's territory.  Thus
these descriptions cover every cyclic order and every assignment of attached
trees.  The weakest resolution is exactly
\[
 \sig(G)>-\delta-\delta+0+0+1
          =1-2\delta=5-2\sqrt5>0.
\]
Therefore all $2116=2110+6$ fully shared $T^6PP$ types are proved without any
rooted hostile-cycle guard.

\section{Main theorem and exhaustion}

\begin{theorem}\label{thm:main}
Every connected octacyclic cactus $G$ on $n$ vertices satisfies
\[
 \boxed{\spE(G)>n}.
\]
\end{theorem}

\begin{proof}
The DNN inequality~\eqref{eq:dnn} proves every cycle multiset except the two in
\eqref{eq:residuals}.  If the shared-cut graph is disconnected, the exact
$44=42+2$ and $76=70+6$ partition audit, the directly verified packing-one
$G7Q$ endpoint, and the strict $877=861+16$ last-bridge certificate prove both
residual families.  If the shared-cut graph is connected, all eight cycles form
one incidence-tree cluster.  The fully shared $T^7Q$ census and common-cut
theorem prove its every incidence; the $2116=2110+6$ census and the six router
resolutions prove every $T^6PP$ incidence.

Every reduced-tree separation is at an actual bridge.  In $G6PP$, the remote
$P_1$ cut leaves all other connector vertices on the marked cluster side.  In
$G7Q$, Lemma~\ref{lem:packing-one} itself permits arbitrary positive path
length and trees on internal path vertices.  Every cycle split uses nonempty
proper consecutive intervals; every shared cut, incidence branch, connector
remnant, and off-hull tree has exactly one owner.  Therefore all territories
are induced, disjoint, and exhaustive, and all packet hypotheses persist.
No edge monotonicity, tree relocation, retracted all-rank guard, or superseded
uncut $868+9$ certificate occurs.  In every case $\sig(G)>0$.
\end{proof}

Since $|E(G)|=n+7$, this proves the strict conclusion of
\cite[Conjecture~1.2]{AKMPZ} for every connected octacyclic cactus.

\appendix
\section{Exact reproduction and dependencies}\label{app:reproduction}

The census and router certificates use the Python standard library and exact
integer or rational arithmetic from Python's \texttt{fractions} module.  The final symbolic
coefficient certificate requires SymPy 1.14.0.  An isolated setup from the
repository root is
\begin{verbatim}
python3 -m venv /tmp/opencode/octacyclic-cacti-venv
/tmp/opencode/octacyclic-cacti-venv/bin/python -m pip install \
  --upgrade pip
/tmp/opencode/octacyclic-cacti-venv/bin/python -m pip install \
  sympy==1.14.0
\end{verbatim}
Then run, in this order,
\begin{verbatim}
/tmp/opencode/octacyclic-cacti-venv/bin/python \
  research/octacyclic-disconnected-partition-census.py
/tmp/opencode/octacyclic-cacti-venv/bin/python \
  research/octacyclic-t6p-last-bridge-conservative.py
/tmp/opencode/octacyclic-cacti-venv/bin/python \
  research/octacyclic-g6pp-last-bridge-census.py
/tmp/opencode/octacyclic-cacti-venv/bin/python \
  research/octacyclic-t6p-last-bridge-sixteen-resolution.py
/tmp/opencode/octacyclic-cacti-venv/bin/python \
  research/octacyclic-g6pp-last-bridge-four-resolution.py
/tmp/opencode/octacyclic-cacti-venv/bin/python \
  research/octacyclic-fully-shared-incidence-census.py
/tmp/opencode/octacyclic-cacti-venv/bin/python \
  research/octacyclic-t6pp-six-exceptions-resolution.py
/tmp/opencode/octacyclic-cacti-venv/bin/python \
  positive-square-energy/experiments/c5_bouquet_matching_certificate.py
\end{verbatim}
The first command asserts $45,44,42,2$ and $77,76,70,6$.  The second and
fourth regenerate $877=861+16$, the profile distribution $(2,5,5,4,0,0)$,
all L1--L16 codes, $16/16$ closure, and weakest margin $1-2\delta$.  The third
and fifth are independent cross-checks.  The sixth enumerates the fully shared
tables and applies the encoded one-cycle SAFE test; it reports the bouquet and
U1--U6 as the unresolved finite incidence types.  The new seventh verifier
regenerates those six rows, checks their stated edge sets, router interval
sizes, sequential branch refinements, retained-cycle packet sets, common-cut
hypotheses, cyclic-cut ownership, and exact symbolic ledgers, and reports
$6/6$ closure.  It imports the census module, so it is independent of neither
the incidence generator nor the SAFE classifier; it does not verify the
analytic packet inequalities.  Likewise, the conservative last-bridge script
imports the marked-census module for finite incidence generation, root orbits,
and encoded packet bounds; the separate sixteen-resolution and independent
G6PP scripts supply ownership and census cross-checks.  None of these scripts
proves induced realization for arbitrary connectors or attached trees:
Lemmas~\ref{lem:whole-cluster} and~\ref{lem:ownership} do that.  The last
command checks only the 1290-term positive-coefficient inequality used in the
two-pentagon common-cut theorem.

Compile from the repository root with
\begin{verbatim}
bash scripts/build-paper.sh all-octacyclic-cacti
\end{verbatim}
or, after installing the packages listed in the accompanying README, with
\begin{verbatim}
latexmk -pdf -interaction=nonstopmode -halt-on-error \
  -cd all-octacyclic-cacti/paper.tex
\end{verbatim}

\begin{thebibliography}{99}
\bibitem{AKMP}
S.~Akbari, H.~Kumar, B.~Mohar, and S.~Pragada,
\emph{A linear lower bound for the square energy of graphs},
Electron. J. Combin. 32(3) (2025), Paper No. P3.53.

\bibitem{AKMPZ}
S.~Akbari, H.~Kumar, B.~Mohar, S.~Pragada, and S.~Zhang,
\emph{Refinement of a conjecture on positive square energy of graphs},
arXiv:2506.07264, 2025.

\bibitem{LTZ}
Y.~Liu, Q.~Tang, and S.~Zhang,
\emph{The positive and negative square-energy conjecture},
arXiv:2607.18031, 2026.

\bibitem{DNN}
Companion manuscript, \emph{The Optimized Liu--Tang--Zhang Constant for
Cactus Graphs}, 2026; repository path \texttt{sharp-cactus-dnn/paper.tex}.

\bibitem{Bicyclic}
Companion manuscript, \emph{Positive Square Energy of Every Bicyclic Cactus},
2026; repository path \texttt{all-bicyclic-cacti/paper.tex}.

\bibitem{Tricyclic}
Companion manuscript, \emph{Positive Square Energy of Every Tricyclic Cactus},
2026; repository path \texttt{all-tricyclic-cacti/paper.tex}.

\bibitem{Tetracyclic}
Companion manuscript, \emph{Positive Square Energy of Every Tetracyclic
Cactus}, repaired strict theorem, 2026; repository path
\texttt{all-tetracyclic-cacti/paper.tex}.

\bibitem{Pentacyclic}
Companion manuscript, \emph{Positive Square Energy of Every Pentacyclic
Cactus}, 2026; repository path \texttt{all-pentacyclic-cacti/paper.tex}.

\bibitem{Hexacyclic}
Companion manuscript, \emph{Positive Square Energy of Every Hexacyclic
Cactus}, 2026; repository path \texttt{all-hexacyclic-cacti/paper.tex}.

\bibitem{Heptacyclic}
Companion manuscript, \emph{Positive Square Energy of Every Heptacyclic
Cactus}, 2026; repository path \texttt{all-heptacyclic-cacti/paper.tex}.

\bibitem{PackingTwo}
Companion manuscript, \emph{Square Energy from Cycle Packing and Block-Graph
Territories}, 2026; repository path
\texttt{packing-two-square-energy/paper.tex}.

\bibitem{Synthesis}
\emph{Positive square energy of every connected octacyclic cactus}, corrected
synthesis, 26 July 2026; repository path
\texttt{research/octacyclic-cactus-complete-synthesis-2026-07-26.md}.

\bibitem{PackingOne}
\emph{Rooted packing-one hostile-cycle lemma for octacyclic provenance},
26 July 2026; repository path
\texttt{research/octacyclic-packing-one-hostile-cycle-lemma-2026-07-26.md}.

\bibitem{LastBridge}
\emph{Conservative last-bridge rerun: all sixteen failed classes},
26 July 2026; repository path
\texttt{research/octacyclic-t6p-last-bridge-conservative-resolution-2026-07-26.md}.

\bibitem{SharedCensus}
\emph{Exact fully shared octacyclic incidence census: $T^7Q$ and $T^6PP$},
26 July 2026; repository path
\texttt{research/octacyclic-fully-shared-incidence-census-2026-07-26.md}.

\bibitem{Six}
\emph{The six fully shared octacyclic $T^6PP$ census exceptions: replacement
audit}, 26 July 2026; repository path
\texttt{research/octacyclic-t6pp-six-exceptions-resolution-2026-07-26.md}.

\bibitem{CommonCut}
\emph{Common-cut residual bouquets: an exact rooted Schur--Sachs inequality},
26 July 2026; repository path
\texttt{research/common-cut-bouquet-rooted-schur-2026-07-26.md}.
\end{thebibliography}
\end{document}
